\begin{explanationbox}

	\textbf{\textcolor{CExplanation}{一句话直觉}}

	圆系方程用一个参数 $\lambda$ 将经过两圆（或圆与直线）公共点的所有圆统一表示，避免逐个求解。

	\medskip
	\textbf{\textcolor{CExplanation}{核心拆解}}
	\begin{description}[style=nextline, leftmargin=2.8em, labelsep=0.5em, itemsep=0.45em]
		\item[\textbf{要点一（线性组合）：}] 将两个方程加权叠加，交点坐标自动满足组合方程，因此曲线必过所有交点。
		\item[\textbf{要点二（参数调节）：}] $\lambda$（或 $\mu$）改变圆的位置和半径；$\lambda=-1$ 时二次项消去，得到公共弦直线（仅两圆系）。
	\end{description}

	\medskip
	\textbf{\textcolor{CExplanation}{几何本质（过定点束）}}

	过一个定点（或点集）的圆构成圆束，圆系方程就是圆束的代数表示。

	\medskip
	\textbf{\textcolor{CExplanation}{代数意义（同解保持）}}

	组合方程是原方程组与参数的线性组合，所有交点坐标仍是新方程的解。

	\medskip
	\textbf{\textcolor{CExplanation}{考点价值（设而不求）}}

	无需直接求交点坐标，设出圆系方程后再代入额外条件确定参数，快速得所求圆方程。

	\medskip
	\textbf{\textcolor{CExplanation}{顿悟点}}

	圆系本质是“过定点的曲线族”，从代数角度统一描述同一公共点集的所有圆。

	\medskip
	\textbf{\textcolor{CExplanation}{使用场景}}
	\begin{itemize}[leftmargin=*, itemsep=0.5em, topsep=0.4em]
		\item \textbf{场景一（过两圆求圆）：}已知两圆相交，再给一个条件（如过某定点），求圆方程时可直接用圆系设而不求。
		\item \textbf{场景二（公共弦方程）：}令圆系中 $\lambda=-1$ 即得公共弦直线。
		\item \textbf{场景三（直线-圆系）：}已知直线与圆相交，类似地设 $C+\mu l=0$，再定 $\mu$。
	\end{itemize}
\end{explanationbox}
